\section{What is object-oriented programming?}
\subsection{Procedural programming and its limitations:}
\begin{itemize}
    \item Procedural programming is pretty much what we've been doing throughout the course.
    \item Focus is on processes or actions that a program takes.
    \item Programs are typically a collection of functions.
    \item Data is declared separately.
    \item Data is passed as arguments to into functions.
    \item Fairly easy to learn.
\end{itemize}
Functions need to know the structure of the data.
\begin{itemize}
    \item If the structure of the data changes many functions must be changed.
    \item Imagine a program with thousands of functions that all depend on a data structure being passed in to it, imagine that that data structure needs to change you would need to change every function, functions might even need to be modified to handle such change and a grave impact is in order, then test all the functions which makes development much more expensive.
\end{itemize}

As programs get larger they become more:
\begin{itemize}
    \item Difficult to understand.
    \item Difficult to maintain.
    \item Difficult to extend.
    \item Difficult to debug.
    \item Difficult to reuse code.
    \item Fragile and easier to break.
\end{itemize}

\subsection{What is objected-oriented programming:}
\begin{itemize}
    \item Object oriented denotes many things such as: OO analysis, OO design, OO testing, etcetera.
    \item Object-oriented is all about modeling you software in terms of objects and classes.
    \item Classes and objects:
        \begin{itemize}
            \item Focus is on classes that model real-world domain entities.
            \item Allows developers to think at a higher level of abstraction.
            \item Used successfully in very large programs.
        \end{itemize}
\end{itemize}

Why is it such a big deal?
\begin{itemize}
    \item As our programs grow more and more complex we need ways to deal with complexity, classes and objects are one way to do just that, rather than thinking of details we can think of just that object.
\end{itemize}

Encapsulation:
\begin{itemize}
    \item Objects contain data AND operations that work on that data.
    \item Abstract Data Type (ADT).
\end{itemize}

Information hiding:
\begin{itemize}
    \item OO allows us to hide implementation specific logic in a class that is available only within the class, now we know that the user can mess with the user implementation specific code since they never knew about it in the first place.
    \item Users of the class code to the interface since they don't need to know the implementation.
    \item More abstraction.
    \item Easier to test, debug, maintain and extend.
\end{itemize}

Reusability:
\begin{itemize}
    \item Easier to reuse in other applications.
    \item Faster development.
    \item Higher quality.
\end{itemize}

Inheritance:
\begin{itemize}
    \item Can create new classes in terms of existing classes.
    \item Reusability.
    \item Polymorphic classes.
\end{itemize}

\subsection{OO limitations}
\begin{itemize}
    \item Not a panacea:
        \begin{itemize}
            \item OO programming won't make bad code better.
            \item Not suitable for all types of problems.
            \item Not everything decomposes to a class.
        \end{itemize}
    
    \item Learning curve:
        \begin{itemize}
            \item Usually a steeper learning curve, especially in C++.
            \item Many OO languages, many variations of OO concepts.
        \end{itemize}
    
    \item Design:
        \begin{itemize}
            \item Usually more up-front design is necessary to create good models and hierarchies.
        \end{itemize}
    
    \item Programs can be:
        \begin{itemize}
            \item Larger in size.
            \item Slower.
            \item More complex.
        \end{itemize}
\end{itemize}


%----------------------------------------------------------------------------------------
\section{What are classes and objects?}
\subsection{Classes}
\begin{itemize}
    \item Blueprints from which objects are created.
    \item A user-defined data-type.
    \item Has attributes (data).
    \item Has methods (functions).
    \item Can hide data and methods.
    \item Provides a public interfase.
\end{itemize}
Example classes:
\begin{itemize}
    \item Account.
    \item Employee.
    \item Image.
    \item \mintinline{cpp}{std::vector}.
    \item \mintinline{cpp}{std::string}.
\end{itemize}

\subsection{Objects}
\begin{itemize}
    \item Objects are created from classes.
    \item Represent a specific instantce of a class.
    \item Can create many, many objects.
    \item Each has its own identity.
    \item Each can use the defined class methods. 
\end{itemize}
Examples of objects:
\begin{itemize}
    \item Frank's account is an instance of an account.
    \item Jim's account is an instance of an account.
    \item Each has its own balance, can make deposits, withdrawls, etc.
\end{itemize}

\subsection{Example}
\begin{itemize}
    \item int are not classes but an object is declared in very much the same way any variable is.
\end{itemize}
\begin{minted}[autogobble]{cpp}
    int high_score;
    int low_score;
    Account frank_account; // user defined type.
    Account jim_account;
    std::vector<int> scores;
    std::string name;
\end{minted}


%----------------------------------------------------------------------------------------
\section{Declaring a class and creating objects}
\begin{minted}[autogobble]{cpp}
    #include <iostream>
    class Player {
        // attributes
        std::string name;
        int health;
        int xp;

        // methods
        void talk(std::string text_to_say);
        bool is_dead();
    };
    int main() {
        // Now that we have declared the class, we can create objects or instances of that class.
        Player *enemy = new Player();
        delete enemy;
        return 0;
    }
\end{minted}

Take the following example:
\begin{minted}[autogobble]{cpp}
    #include <iostream>
    class Account {
        // attributes
        std::string name;
        double balance;
        // methods
        bool withdrawl(double amount);
        bool deposit(double amount);
    };
    int main() {
        Account frank_account;
        Account jim_account;
        Account *mary_account = new Account();
        delete mary_account;
        return 0;
    }
\end{minted}

\subsection{Creating objects}
\begin{itemize}
    \item Once we have objects we can use it like any other variable, for example.
        \begin{minted}[autogobble]{cpp}
            Account frank_account, jim_account;
            std::vector<Account> account1 {frank_account, jim_account};
            account1.push_back(jim_account);
            Account accounts[] {frank_account};
        \end{minted}
\end{itemize}

\subsection{Example}
\begin{minted}[autogobble]{cpp}
    #include <iostream>
    #include <string>
    #include <vector>
    using namespace std;
    class Player { // everything inside here is 'encapsulated'.
        // attributes 
        string name;
        unsigned int health;
        int xp;
        // methods prototypes
        void talk(string);
        bool is_dead();
    };
    class Account {
        string name;
        double balance;
        
        bool deposit(double);
        bool withdraw(double);
    };
    int main() {
        Player frank, hero;
        Player * enemy {nullptr};
        enemy = new Player;
        delete enemy;

        Player players[] {frank, hero};

        vector<Player> player_vec;
        player_vec.push_back(hero);

        Account frank_account, jim_account;
        return 0;
    }
\end{minted}

%----------------------------------------------------------------------------------------
\section{Accessing class members}
\begin{itemize}
    \item We can access:
        \begin{itemize}
            \item Class attributes.
            \item Class methods.
        \end{itemize}
    
    \item Some class members will not be accessible (because of the access modifiers).
    \item We need an object to access instance variables.
\end{itemize}

\subsection{Access members}
\begin{itemize}
    \item If we have an object then we can use the dot operator to access the members.
        \begin{minted}[autogobble]{cpp}
            frank_account.balance;
            frank_account.deposit(1000.00);
        \end{minted}
    
    \item If we have a pointer to an object then we can use the \mintinline{cpp}{->} operator or dereference and use the dot operator to access the members.
        \begin{minted}[autogobble]{cpp}
            Account *frank_account = new Account();
        (*frank_account).balance;
        (*frank_account).deposit(1000.00);
        \end{minted}
        \begin{minted}[autogobble]{cpp}
            Account *frank_account = new Account();
            frank_account->balance;
            frank_account->deposit(1000.00);
        \end{minted}
\end{itemize}

\subsection{Example}
\begin{minted}[autogobble]{cpp}
    #include <iostream>
    using namespace std;

    class Player {
        public: // allows us to see the attributes and methods from main.
        string name;
        double health;
        unsigned int xp;

        void talk(string text_to_say) {
            cout << name << " says " << text_to_say << endl;
        }
        bool is_dead(); // not implemented, if executed a linker error is going to be thrown.
    };

    int main() {
        Player frank;
        frank.name = "Frank"; // access the name and change it to "Frank"
        frank.health = 100; // acces the health and change it to 100.
        frank.xp = 12; // access the xp and change it 12.
        frank.talk("Hello");

        Player *enemy = new Player;
        (*enemy).name = "Enemy";
        enemy->name = "Enemy";
        (*enemy).health = 100;
        enemy->health = 100;
        (*enemy).xp = 12;
        enemy->xp = 12;
        (*enemy).talk("I will destroy you.");
        enemy->talk("I will destroy you.");
        
        return 0;
    }
    /* OUTPUT:
    Frank says Hello
    Enemy says I will destroy you.
    Enemy says I will destroy you.
    */
\end{minted}


%----------------------------------------------------------------------------------------
\section{Public and private}
\begin{itemize}
    \item \mintinline{cpp}{public:}
        \begin{itemize}
            \item Accessible everywhere.
        \end{itemize}
    
    \item \mintinline{cpp}{private:}
        \begin{itemize}
            \item Accessible only by memebers or friends (we don't know yet) of the class.
        \end{itemize}
    
    \item \mintinline{cpp}{protected:}
        \begin{itemize}
            \item Used with inheritance (we don't know that yet.)
        \end{itemize}
\end{itemize}

\subsection{Declaration}
\begin{itemize}
    \item Once you declare an access modifier, everything from that point forward will have tha access modifier applied to it.
        \begin{minted}[autogobble]{cpp}
            class Class_name {
                public: 
                // declaration(s);
            };
        \end{minted}
    
    \item By default everything in a class will be private if no access modifier is specified.
        \begin{minted}[autogobble]{cpp}
            class Class_name {
                private: 
                // declaration(s);
            };
        \end{minted}
    
    \item Protected comes in when using inheritance, its basically the private access modifier for inheritance.
        \begin{minted}[autogobble]{cpp}
            class Class_name {
                protected:
                // declaration(s);
            };
        \end{minted}
    
    \item It is common to combine access modifiers because we want to hide only certain things.
        \begin{minted}[autogobble]{cpp}
            class Player {
                private:
                std::string name;
                unsigned int health;
                int xp;
                public:
                void talk(std::string text_to_say);
                bool is_dead();
            };
        \end{minted}
    
    \item If you try to access members declared private you will get a compiler error.
        \begin{minted}[autogobble]{cpp}
            #include <iostream>
            using namespace std;

            class Player {
                private:
                std::string name;
                unsigned int health;
                int xp;
                public:
                void talk(std::string text_to_say);
                bool is_dead();
            };


            int main() {
                Player frank;
                frank.name = "Frank"; // compiler error.
                frank.health = 100; // compiler error.
                frank.talk("Ready to battle.");
                return 0;
            }
        \end{minted}
\end{itemize}


%----------------------------------------------------------------------------------------
\section{Implementing member methods}
\begin{itemize}
    \item Very similar to how we implemented functions.
    \item Member methods have access to the member attributes.
        \begin{itemize}
            \item So you don't need to pass them as arguments.
        \end{itemize}

    \item Can be implemented outside the class declarationg.
        \begin{itemize}
            \item Need to use \verb|class_name::method_name|.
        \end{itemize}
    
    \item Can separate specification from implementation.
        \begin{itemize}
            \item .h file for the class declaration.
            \item .cpp file for the class implementation.
        \end{itemize}
\end{itemize}

\subsection{Inside the class declaration}
\begin{minted}[autogobble]{cpp}
    class Account {
        private:
        double balance;
        public:
        void set_balance(double bal) {
            balance = bal;
        }
        double get_balance() {
            return balance;
        }
    };
\end{minted}

\subsection{Outside the class declaration}
Just have the method prototypes inside the class.
\begin{minted}[autogobble]{cpp}
    class Account {
        private:
        double balance;
        public:
        void set_balance(double bal);
        double get_balance();
    };

    void Account::set_balance(double bal) {
        balance = bal;
    }

    void Account::get_balance() {
        return balance;
    }
\end{minted}

\subsection{Separate specification from implementation}
An include-guard is included so that if you include the file more than once you won't have a class redefinition error.
\begin{minted}[autogobble]{cpp}
    #ifndef _ACCOUNT_H_ // if this is not defined
    #define _ACCOUNT_H_ // define it. Otherwise skip to #endif

    #endif
\end{minted}
Some compilers allow you to include-guard your files with a simple preprocessor directive called \mintinline{cpp}{#pragma}.
\begin{minted}[autogobble]{cpp}
    #pragma once 
    class Account {};

\end{minted}
Specify the class in a file ``Account.h'':
\begin{minted}[autogobble]{cpp}
    #ifndef _ACCOUNT_H_
    #define _ACCOUNT_H_

    class Account {
        private:
        double balance;
        public:
        void set_balance(double bal);
        double get_balance();
    };

    #endif
\end{minted}
Implement the class in ``Account.cpp''
\begin{minted}[autogobble]{cpp}
    #include "Account.h"
    void Account::set_balance(double bal) {
        balance = bal;
    }
    double Account::get_balance() {
        return balance;
    }
\end{minted}

\subsection{Example}
\begin{minted}[autogobble]{cpp}
    #include <iostream>
    using namespace std;

    class Account {
    private:
        string name;
        double balance;
    public:
        void set_balance(double bal) {balance = bal;}
        double get_balance() {return balance;}

        void set_name(string n);
        string get_name();
        bool deposit(double amount);
        bool withdraw(double amount);
    };

    void Account::set_name(string n) {
        name = n;
    }
    string Account::get_name() {
        return name;
    }
    bool Account::deposit(double amount) {
        balance += amount;
        return true;
    }
    bool Account::withdraw(double amount) {
        if (balance-amount >= 0) {
            balance -= amount;
            return true;
        }
        else {
            return false;
        }
    }


    int main() {
        Account frank_account;
        frank_account.set_name("Frank's account");
        frank_account.set_balance(1000.00);
        if (frank_account.deposit(200.00)) {
            cout << "Deposit OK" << endl;
        }
        else {
            cout << "Deposit Not allowed." << endl;
        }

        if (frank_account.withdraw(500.00)) {
            cout << "Withdraw OK" << endl;
        } 
        else {
            cout << "Not sufficient funds." << endl;
        }

        if (frank_account.withdraw(1500.00)) {
            cout << "Withdraw OK" << endl;
        }
        else {
            cout << "Not sufficient funds" << endl;
        }
        return 0;
    }
    /* OUTPUT:
    Deposit OK
    Withdraw OK
    Not sufficient funds
    */
\end{minted}



%----------------------------------------------------------------------------------------
\section{Constructors and destructors}
\begin{itemize}
    \item Constructors are special member methods that are invoked during object creation, they are commonly used for initialization, constructors have the same name as the class and can be overloaded. 
    
    \item Destructors are also special member methods, they bear the same name as the class proceeded with a tilde \verb|~|, destructors are invoked automatically when an object is destroyed, they have no return type and no parameters, only 1 destructor is alowed per class and these cannot be overloaded. They are useful to release memory, close files, and other resources. You can call the destructor using \mintinline{cpp}{delete}, however if you don't call it, once the object goes out of scope it will automatically be deleted by C++.
    \item Example of constructors and destructors:
    \begin{minted}[autogobble]{cpp}
        class Player {
            private:
                std::string name;
                int health;
                int xp;
            public:
                // Constructors.
                Player();
                Player(std::string name);
                Player(std::string name, int health, int xp);
                // Destructors.
                ~Player();
        };

        class Account {
            private:
                std::string name;
                double balance;
            public: 
                // Constructors.
                Account();
                Account(std::string name, double balance);
                Account(std::string name);
                Account(double balance);
                // Destructors.
                ~Account();
        };
    \end{minted}
\end{itemize}

\subsection{Creating objects}
\begin{itemize}
    \item Objects can be called with any type of initialization and depending on the types of the parameters C++ will determine which one to use.
        \begin{minted}[autogobble]{cpp}
            #include <iostream>
            using namespace std;
            class Player {
                private:
                    std::string name;
                    int health;
                    int xp;
                public:
                // Constructors
                Player() {
                    cout << "No args constructor called." << endl;
                }
                Player(std::string name) {
                    cout << "String arg constructor called." << endl;
                }
                Player(std::string name, int health, int xp) {
                    cout << "Three args constructor called" << endl;
                }
                // Destructor.
                ~Player() {
                    cout << "Destructor called for " << name << endl;
                }
                void set_name(std::string name) {
                    this->name = name;
                }
            };

            class Account {
                private:
                    std::string name;
                    double balance;
                public: 
                    // Constructors.
                    Account();
                    Account(std::string name, double balance);
                    Account(std::string name);
                    Account(double balance);
            };

            int main() {
                {
                    Player slayer;
                    slayer.set_name("Slayer");
                }
                /* OUTPUT:
                No args constructor called.
                Destructor called for Slayer
                */
                {
                    Player frank;
                    frank.set_name("Frank");
                    Player hero("Hero");
                    hero.set_name("Hero");
                    Player villain("Villain", 100, 12);
                    villain.set_name("Villain");
                }
                /* OUTPUT: Notice the objects are destroyed in reverse order because they are in a stack.
                No args constructor called.
                String arg constructor called.
                Three args constructor called
                Destructor called for Villain
                Destructor called for Hero
                Destructor called for Frank
                */
                {
                    Player *enemy = new Player; // No args constructor.
                    enemy->set_name("Enemy");
                    Player *level_boss = new Player("Level Boss", 1000, 300); // 3 args constructor called.
                    level_boss->set_name("Level Boss");
                    delete enemy; // destructor called.
                    delete level_boss; // destructor called.
                }
                /* OUTPUT:
                No args constructor called.
                Three args constructor called
                Destructor called for Enemy
                Destructor called for Level Boss
                */
                return 0;
            }
        \end{minted}
\end{itemize}


%----------------------------------------------------------------------------------------
\section{The default constructor}
\begin{itemize}
    \item A default constructor does not expect any arguments, it is also called a no-args constructor. If you write no constructor at all for the class then C++ will generate a default constructor that does nothing. This default constructor is the one that is called when you create or instantiate a new object with no arguments.
    \item We are always free to provide our no args constructor and in fact it is best practice to do so.
    \item When there is at least one constructor in the class, C++ will not generate a dummy constructor, if we still need the no args constructor we must define it ourselves.
\end{itemize}

\subsection{Example}
\begin{minted}[autogobble]{cpp}
    #include <iostream>
    using namespace std;

    class Player {
        // C++ will generate a dummy constructor behind the scenes since there is none defined.
        private:
            std::string name;
            int health;
            int xp;
        public:
        void set_name(std::string name_val) {
            name = name_val;
        }
        std::string get_name() {
            return name;
        };
    };

    class Player_ {
        // C++ will not generate a dummy constructor, we have provided at least one.
        private:
            std::string name;
            int health;
            int xp;
        public:
        Player_() { // this is the default no args initializer.
            name = "None";
            health = 100;
            xp = 3;
        }
        void set_name(std::string name_val) {
            name = name_val;
        }
        std::string get_name() {
            return name;
        };
    };

    int main() {
        Player frank;
        frank.set_name("Frank");
        cout << frank.get_name() << endl;
        Player_ frank_with_constructor;
        cout << frank_with_constructor.get_name() << endl;
        return 0;
    }
    /* OUTPUT:
    Frank
    None
    */
\end{minted}


%----------------------------------------------------------------------------------------
\section{Overloading constructor}
\begin{itemize}
    \item Classes can have as many constructors as necessary, each must have a unique signature. The compiler has to be able to determine which to call based on the initialization information provided when creating the objects, if there is any ambiguity, the compiler won't guess, it will instead generate a compiler error. Default constructors are no longer provided for you classes once you've declared at least one.
\end{itemize}

\subsection{Example}
\begin{minted}[autogobble]{cpp}
    #include <iostream>
    using namespace std;

    class Player {
        private:
        std::string name;
        int health, xp;
        public:
        Player();
        Player(std::string name_val);
        Player(std::string name_val, int health_val, int xp_val);
    };

    Player::Player() {
        name = "None";
        health = 0;
        xp = 0;
    }

    Player::Player(std::string name_val) {
        name = name_val;
        health = 0;
        xp = 0;
    }

    Player::Player(std::string name_val, int health_val, int xp_val) {
        name = name_val;
        health = health_val;
        xp = xp_val;
    }

    int main() {
        Player empty; // None, 0, 0
        Player hero {"Hero"}; // Hero, 0, 0
        Player villain {"Villain"}; // Villain, 0, 0
        Player frank {"Frank", 100, 4}; // Frank, 100, 4
        Player *enemy = new Player("Enemy", 1000, 0); // Enemy, 1000, 0
        delete enemy;
        return 0;
    }
\end{minted}


%----------------------------------------------------------------------------------------
\section{Constructor Initialization lists} 
\begin{itemize}
    \item So far, all data memeber values have been set in the constructor body. What we really want to do is have the member data values initialized to our values before the constructor body executes, this is more efficient.
    \item Constructor initialization lists are: 
        \begin{itemize}
            \item More efficient.
            \item Initialization list immediately follows the parameter list.
            \item Initializes the data members as the objects created.
            \item Order of initialization is the order of declaration in the class.
        \end{itemize}
    
    \item Take the following example that does NOT use constructor initialization lists:
        \begin{minted}[autogobble]{cpp}
            Player::Player() {
                name = "None";
                health = 0;
                xp = 0;
            }
        \end{minted}
        \begin{itemize}
            \item While this works it takes more time and is inefficient. A better way is to provide a constructor initialization list.
            \item By the time we get to the statement \mintinline{cpp}{name = "None";} the string object has already been constructed and initialized to an empty string, in other words \mintinline{cpp}{name = "None";} just assigns to an already initialized object and doesn't initialize the object.
            \item So in order to initialize the members as they are created as to have a true initialization we must use a constructor initialization list.
        \end{itemize}
    
    \item Take the following example that does use constructor initialization lists:
        \begin{minted}[autogobble]{cpp}
            Player::Player() : name{"None"}, health{0}, xp{0} {
            }
        \end{minted}
        \begin{itemize}
            \item The above actually initializes before doing anything, this initializes everything before the body of the constructor is ever excecuted.
            \item The data memebers will not be initialized exactly in the order provided in the constructor initialization list, they will be initialized in the order in which they are provided in the class declaration.
        \end{itemize}
\end{itemize}

\subsection{Changing previous initialization to use constructor initialization lists}
\begin{minted}[autogobble]{cpp}
    #include <iostream>
    using namespace std;

    class Player {
        private:
        std::string name;
        int health, xp;
        public:
        Player();
        Player(std::string name_val);
        Player(std::string name_val, int health_val, int xp_val);
    };

    Player::Player() : name{"None"}, health{0}, xp{0} {
    }

    Player::Player(std::string name_val) : name{name_val}, health{0}, xp{0} {
    }

    Player::Player(std::string name_val, int health_val, int xp_val) : name{name_val}, health{health_val}, xp{xp_val}  {
    }

    int main() {
        Player empty; // None, 0, 0
        Player hero {"Hero"}; // Hero, 0, 0
        Player villain {"Villain"}; // Villain, 0, 0
        Player frank {"Frank", 100, 4}; // Frank, 100, 4
        Player *enemy = new Player("Enemy", 1000, 0); // Enemy, 1000, 0
        delete enemy;
        return 0;
    }
\end{minted}


%----------------------------------------------------------------------------------------
\section{Delegating constructors}
\begin{itemize}
    \item Often the code for constructors is very similar, duplicated code can lead to errors. C++ allows delegating constructors which is code for one constructor can call another in the initialization list, this in turn avoids duplicating code.
    \item This only works it constructors are called using constructor initialization lists, you cannot call constructors explicitly from the body of another. 
    \item When you call other constructors, the body of the called constructors will excecute and then excecution will continue from the constructor.
    \item So remember, if you delegate the constructors, both bodies of the constructors will be excecuted, the called constructor's body will be excecuted first and then the rest of the constructor body is excecuted.
\end{itemize}

\subsection{Example}
\begin{minted}[autogobble]{cpp}
    #include <iostream>
    using namespace std;

    class Player {
        private:
        std::string name;
        int health, xp;
        public:
        Player();
        Player(std::string name_val);
        Player(std::string name_val, int health_val, int xp_val);
    };

    Player::Player() : Player {"None", 0, 0 } {
        cout << "No args constructor" << endl;
    }

    Player::Player(std::string name_val) : Player{name_val, 0, 0} {
        cout << "One arg constructor" << endl;
    }

    Player::Player(std::string name_val, int health_val, int xp_val) : name{name_val}, health{health_val}, xp{xp_val}  {
        cout << "Three args constructor" << endl;
    }

    int main() {
        Player empty; 
        Player hero {"Hero"}; 
        Player frank {"Villain", 100, 55}; 
        return 0;
    }
    /* OUTPUT:
    Three args constructor
    No args constructor
    Three args constructor
    One arg constructor
    Three args constructor
    */
\end{minted}


%----------------------------------------------------------------------------------------
\section{Constructor parameters and default values}
\begin{itemize}
    \item We can use default values to simplify our code and reduce the number of overloaded constructors.
    \item Same rules apply as we learned with non-member functions.
    \item When doing this it is important to provide unambiguous constructors, for example if we declare one constructor with all default values and then declare a no arg constructor we would get a compiler error because C++ would not know which one to call.
\end{itemize}

\subsection{Example}
\begin{minted}[autogobble]{cpp}
    #include <iostream>
    using namespace std;

    class Player {
        private:
        std::string name;
        int health, xp;
        public:
        Player(std::string name_val = "None", int health_val = 0, int xp_val = 0);
    };

    Player::Player(std::string name_val, int health_val, int xp_val) : name{name_val}, health{health_val}, xp{xp_val}  {
        cout << "Three args constructor" << endl;
    }

    int main() {
        Player empty; 
        Player frank {"Frank"};
        Player hero {"Hero", 100};
        Player villain {"Villain", 100, 55};
        return 0;
    }
    /* OUTPUT:
    Three args constructor
    Three args constructor
    Three args constructor
    Three args constructor
    */
\end{minted}


%----------------------------------------------------------------------------------------
\section{Copy constructor}
\begin{itemize}
    \item When objects are copied C++ must create a new object from an existing object.
    \item When is a copy of an object made?
        \begin{itemize}
            \item Passing object by value as parameter.
            \item Returning an object from a function by value.
            \item Constructing one object based on another of the same class.
        \end{itemize}
    
    \item C++ must have a way of accomplish this copying, so C++ uses a copy constructor. C++ will provide a compiler defined copy constructor if you don't.
\end{itemize}

\subsection{1st usecase: Pass object by-value}
Example:
\begin{minted}[autogobble]{cpp}
    Player hero {"Hero", 100, 20};
    void display_player(Player p) {
        // p is a copy of hero in this example.
        // use p
        // Destructor for p will be called.
    }
    display_player(hero);
\end{minted}

\subsection{2nd usecase: Return object by value}
Example:
\begin{minted}[autogobble]{cpp}
    Player enemy;
    Player create_super_enemy() {
        Player an_enemy {"Super Enemy", 1000, 1000};
        return an_enemy; // A copy of an_enemy is returned, the copy is made by the copy constructor.
    }
    enemy = create_super_enemy();
\end{minted}

\subsection{Construct one object based on another}
Example:
\begin{minted}[autogobble]{cpp}
    Player hero{"Hero", 100, 100};
    Player another_hero{hero}; // A copy of hero is made.
\end{minted}

\subsection{Copy constructors}
\begin{itemize}
    \item If you don't provide your own way of copying objects by value then the compiler provides a default way of copying objects.
    \item The copy constructors copy the values of ach data memeber to the new object, this is called 'default-memberwise-copy' this basically means that for each data member C++ will copy them to the new object.
    \item In many cases this is fine, except for one thing, pointers.
    \item Beware if you have a pointer data member, because the pointer will be copied but data that the pointer is pointing to will not. This is explained in the next video.
\end{itemize}

\subsection{Best practice is:}
\begin{itemize}
    \item Provide a copy constructor when your class has raw pointer members.
    \item Provide the copy constructor with a const reference parameter.
    \item Use STL classes as they already provide copy constructors.
    \item Avoid using raw pointer data members if possible.
\end{itemize}

\subsection{Declaring the copy constructor}
\begin{itemize}
    \item Copy constructors are declared like this:
        \begin{minted}[autogobble]{cpp}
            Type::Type(const type &source);
            Player::Player(const Player &source);
            Account::Account(const Account &source);
        \end{minted}
    
    \item Actually implementing the copy constructor is like this:
        \begin{minted}[autogobble]{cpp}
                Player::Player(const Player &source) : name{source.name}, health{source.health}, xp{source.xp} {
                }
                Account::Account(const Account &source) : name{source.name}, balance{source.balance} {
                }
        \end{minted}
\end{itemize}

\subsection{Example}
\begin{minted}[autogobble]{cpp}
    #include <iostream>
    using namespace std;

    class Player {
        private:
        std::string name;
        int health;
        int xp;
        public: 
        std::string get_name() {return name;}
        int get_health() {return health;}
        int get_xp() {return xp;}
        Player(std::string name_val = "None", int health_val = 0, int xp_val = 0);
        // Copy constructor:
        Player(const Player &source);
        // Destructor:
        ~Player() { cout << "Destrocutor called for: " << name << endl; }
    };

    Player::Player(std::string name_val, int health_val, int xp_val) : name{name_val}, health{health_val}, xp{xp_val} {
        cout << "Theree-args constructor for " + name << endl;
    }

    Player::Player(const Player &source) : name{source.name}, health{source.health}, xp{source.xp} {
        cout << "Copy constructor - made a copy of: " << source.name << endl;
    }

    void display_player(Player p) {
        cout << "Name: " << p.get_name() << endl;
        cout << "Health: " << p.get_health() << endl;
        cout << "xp: " << p.get_xp() << endl;
    }

    int main() {
        Player empty;
        display_player(empty);
        Player copy_of_empty {empty};
        display_player(copy_of_empty);
        Player frank {"Frank"};
        Player hero {"Hero", 100};
        Player villain {"Villain", 100, 55};
        return 0;
    }
    /* OUTPUT:
    Theree-args constructor for None
    Copy constructor - made a copy of: None
    Name: None
    Health: 0
    xp: 0
    Destrocutor called for: None
    Copy constructor - made a copy of: None
    Copy constructor - made a copy of: None
    Name: None
    Health: 0
    xp: 0
    Destrocutor called for: None
    Theree-args constructor for Frank
    Theree-args constructor for Hero
    Theree-args constructor for Villain
    Destrocutor called for: Villain
    Destrocutor called for: Hero
    Destrocutor called for: Frank
    Destrocutor called for: None
    Destrocutor called for: None
    */
\end{minted}

